% Original Bengali explanations after comparison; sources in CHAPTER04_BOOK_COMPARISON.md.
\subsubsection{হাইপারটেক্সট ও হাইপারমিডিয়া}
\textbf{হাইপারটেক্সট (Hypertext)} হলো লিংকের মাধ্যমে সম্পর্কিত নথি বা নথির অংশে যাওয়ার উপযোগী পাঠ। বইয়ের সূচি থেকে অধ্যায়ে যাওয়ার মতো, তবে ওয়েবে লিংক নির্বাচন করেই গন্তব্যে পৌঁছানো যায়। \textbf{হাইপারমিডিয়া (Hypermedia)} এই সংযোগের ধারণাকে লেখা ছাড়াও ছবি, শব্দ ও ভিডিওর মতো মাধ্যমে প্রসারিত করে। যেমন, বিজ্ঞান ক্লাবের পেজের একটি যন্ত্রের ছবি থেকে তার কাজের ব্যাখ্যা বা প্রদর্শনের ভিডিওতে যাওয়া যায়। শুধু একই পেজে নানা মাধ্যম থাকলে সেটি multimedia; লিংক দিয়ে সম্পর্কিত বিষয় অনুসরণের দিকটি hypermedia-তে গুরুত্বপূর্ণ।

\subsubsection{ওয়েব পোর্টাল: সেবার প্রবেশদ্বার}
বিভিন্ন তথ্য, সেবা ও সংশ্লিষ্ট সাইটের লিংক এক জায়গায় সাজিয়ে দেওয়া বিশেষ ধরনের ওয়েবসাইটকে \textbf{ওয়েব পোর্টাল (Web Portal)} বলা হয়। একটি শিক্ষার্থী পোর্টালে নোটিশ, কোর্স, পরীক্ষার রুটিন ও ফলাফলের প্রবেশপথ থাকতে পারে। প্রতিটি সেবার ফাইল একই সার্ভারে থাকতেই হবে না। \textbf{পার্থক্য:} homepage একটি সাইটের প্রধান পেজ; portal একটি সাইটের উদ্দেশ্য বা সেবার ধরন। তাই প্রতিটি homepage-কে portal বলা ঠিক নয়।

কলেজের রুটিন, ফি ও ফলাফলের সেবায় এক জায়গা থেকে যাওয়ার ব্যবস্থা থাকলে সেই সাইট পোর্টালের উদাহরণ। শুধু সাধারণ পরিচিতির পেজে এই বৈশিষ্ট্য নাও থাকতে পারে।
